\documentclass[10pt,a4paper]{amsart}
\usepackage[margin=19mm]{geometry}
\usepackage{amsmath,amssymb,mathtools}
\usepackage[scr=rsfs,cal=euler]{mathalfa}
\usepackage{xfrac}
\usepackage[unicode,hidelinks,bookmarks=false]{hyperref}
\newcommand{\ie}{i.e.\ }
\newcommand{\cf}{cf.\ }
\newcommand{\resp}{respectively\ }
\newcommand{\Subsection}{\subsection}
\providecommand{\nobreakdash}{\nobreak\hbox{-}}
\setlength{\emergencystretch}{2em}
\multlinegap0pt
\usepackage{amssymb}
\usepackage{enumerate}
\usepackage{tikz}
\newtheorem{lem}{Lemma}
\newcommand{\bu}{\mathbf{u}}
\title{Generalized $q$-dimensions of measures on Non-autonomous conformal sets}
\author{Jun Jie Miao, Tianrui Wang}
\date{Source: arXiv:2512.19771v1 (CC BY 4.0)}
\begin{document}
\maketitle

\noindent\emph{Source and licence.} Generalized $q$-dimensions of measures on Non-autonomous conformal sets. Version arXiv:2512.19771v1 (CC BY 4.0), distributed by arXiv under the Creative Commons Attribution 4.0 International licence, http://creativecommons.org/licenses/by/4.0/, as recorded in the arXiv licence metadata for this exact version and shown on the abstract page https://arxiv.org/abs/2512.19771v1. Full licence notice: \texttt{licenses/arxiv-2512.19771-CC-BY-4.0.txt}.\par\smallskip

\noindent\textit{Editorial scope.} Counted unit: the article's lem d (source0.tex lines 477--481). The article's complete original proof of it is delivered with it (source0.tex lines 483--534, one \texttt{proof} environment of 52 lines). No mathematics is written, repaired or re-derived: the statement and the proof are the article's own lines, in the article's own notation, and no retained line is substituted or re-flowed.  Retained uncounted as the source-local context the proof needs: lines 150--152; lines 248--257; lines 260--265; lines 455--461.  Nothing was omitted from any retained span, so the package carries no omission record.  No retained line contains a citation, so the wrapper carries no bibliography and no bibliography entry.  No retained line contains a figure environment, an image-inclusion command or drawing code, so no image is delivered and the package declares no asset dependency.  Editorial formatting only, in addition: an AMS \texttt{amsart} preamble that restates the article's own declarations and macros used by the retained lines, namely the environments \texttt{lem} on separate counters, together with the standard packages those lines need.  The retained environments print in this document as Lemma 1 in order of appearance, whereas the article prints the same items with the numbering of its own full document.  The complete native source of the article as distributed in the arXiv e-print for this exact version is bundled unchanged as \texttt{sources/191507/source0.tex}.
\begin{equation} \label{def_uccdn}
|\varphi_{k, i}(x)-\varphi_{k, i}(y)|\le  c|x-y| \qquad \textit{ for all }  x, y\in J.
\end{equation}

\begin{equation}\label{def_GULPF}
\begin{split}
&\overline P_\mu(t, q)=\left\{ \begin{aligned} &\limsup_{\delta\to 0} \frac{\mathrm{sgn}(1-q)}{k_{\delta}}\log \sum_{\mathbf{u}\in\mathcal{C}(\delta)}\|D\varphi_\mathbf{u}\|^{t(1-q)}\mu([\bu])^q ,  \textit{  $q>0$, $q\not=1$}   \\
&\limsup_{\delta\to 0} -\frac{1}{k_{\delta}} \sum_{\mathbf{u}\in\mathcal{C}(\delta)}\mu([\bu])\log(\|D\varphi_\mathbf{u}\|^{-t}\mu([\bu])), \hspace{1cm}q=1,
\end{aligned} \right.   \\
&\underline P_\mu(t, q)=\left\{ \begin{aligned} & \liminf_{\delta\to 0} \frac{\mathrm{sgn}(1-q)}{k_{\delta}}\log \sum_{\mathbf{u}\in\mathcal{C}(\delta)}\|D\varphi_\mathbf{u}\|^{t(1-q)}\mu([\bu])^q ,  \textit{  $q>0$, $q\not=1$}   \\
&\liminf_{\delta\to 0} -\frac{1}{k_{\delta}} \sum_{\mathbf{u}\in\mathcal{C}(\delta)}\mu([\bu])\log(\|D\varphi_\mathbf{u}\|^{-t}\mu([\bu])), \hspace{1cm} q=1.
\end{aligned} \right.
\end{split}
\end{equation}    

\begin{equation}\label{d*}
\begin{split}
\overline{d}_q^*=\inf\{t:\overline P_\mu(t, q)<0\}=\sup\{t:\overline P_\mu(t, q)>0\},  \\
\underline{d}_q^*=\inf\{t:\underline P_\mu(t, q)<0\}=\sup\{t:\underline P_\mu(t, q)>0\}.
\end{split}
\end{equation}             

\begin{lem}\label{u contraction}
Given a $NCIFS$, for each integer $k$ and $u\in I_k$, 
$$
\|D\varphi_u\|\le c,
$$
where $c$ is given by \eqref{def_uccdn}.
\end{lem}

\begin{lem}\label{m d}
Both  $\overline P_\mu(t, q)$ and  $\underline P_\mu(t, q)$  in \eqref{def_GULPF} are  monotonically  decreasing  in $t$. 

In particular, given $q>0$,   if $\overline P_\mu(t, q)$ and  $\underline P_\mu(t, q)$ is finite on an  interval $I$, then they are strictly decreasing on $I$,  convex  when $0<q<1$ and  concave   when $q>1$. Moreover $\overline{d}_q^*$ and $\underline{d}_q^*$ in \eqref{d*} are finite. 
\end{lem}

\begin{proof}
Given $q\in(0, 1)$,  suppose $\overline P_\mu(t_1, q)$ and $\overline P_\mu(t_2, q)$ are finite for $t_2>t_1$. It follows by Lemma \ref{u contraction} that
$$
\sum_{\mathbf{u}\in\mathcal{C}(\delta)}\|D\varphi_\mathbf{u}\|^{t_2(1-q)}\mu([\bu])^q
\le c^{k_\delta(1-q)(t_2-t_1)}\sum_{\mathbf{u}\in\mathcal{C}(\delta)}\|D\varphi_\mathbf{u}\|^{t_1(1-q)}\mu([\bu])^q. 
$$
Since $0<c<1$, by \eqref{def_GULPF},  we have that
$$
\overline P_\mu(t_2, q)\le\overline P_\mu(t_1, q),   \qquad  \underline P_\mu(t_2, q)\le\underline P_\mu(t_1, q).
$$
For $q=1$ and $q>1$,  by the similar argument,  we have  $\overline P_\mu(t_2, q) \leq  \overline P_\mu(t_1, q)$ and $\underline P_\mu(t_2, q)\le\underline P_\mu(t_1, q)$.


\iffalse
$\overline P_\mu(t_2, q) < \overline P_\mu(t_1, q)$.

Let $q=1$. By \eqref{def_uccdn}, we have
$$
\sum_{\mathbf{u}\in\mathcal{C}(\delta)}\mu([\bu])\log(\|D\varphi_\mathbf{u}\|^{t_2}\mu([\bu])^{-1})\le k_\delta(t_2-t_1)\log c+\sum_{\mathbf{u}\in\mathcal{C}(\delta)}\mu([\bu])\log(\|D\varphi_\mathbf{u}\|^{t_1}\mu([\bu])^{-1}).
$$
Thus
$$
\overline P_\mu(t_2, q)\le\overline P_\mu(t_1, q)-(t_2-t_1)\log\frac{1}{c}.
$$
Since $0<c<1$, $\overline P_\mu(t_2, q) < \overline P_\mu(t_1, q)$.

Let $q>1$. By \eqref{def_uccdn}, we have
$$
\sum_{\mathbf{u}\in\mathcal{C}(\delta)}\|D\varphi_\mathbf{u}\|^{t_2(1-q)}\mu([\bu])^q \ge c^{k_\delta(1-q)(t_2-t_1)}\sum_{\mathbf{u}\in\mathcal{C}(\delta)}\|D\varphi_\mathbf{u}\|^{t_1(1-q)}\mu([\bu])^q.
$$
Therefore
$$
\overline P_\mu(t_2, q)\le\overline P_\mu(t_1, q)-(1-q)(t_2-t_1)\log\frac{1}{c}.
$$
Since $0<c<1$, $\overline P_\mu(t_2, q) < \overline P_\mu(t_1, q)$.
\fi


For $q>0, q\not=1$ and $0<\alpha<1$, by Hölder inequality, we have
\begin{eqnarray*}
&&\sum_{\mathbf{u}\in\mathcal{C}(\delta)}\|D\varphi_\mathbf{u}\|^{(\alpha t_1+(1-\alpha) t_2)(1-q)}\mu([\bu])^q   \\
&=&\sum_{\mathbf{u}\in\mathcal{C}(\delta)}(\|D\varphi_\mathbf{u}\|^{t_1(1-q)}\mu([\bu])^q)^\alpha(\|D\varphi_\mathbf{u}\|^{t_2(1-q)}\mu([\bu])^q)^{1-\alpha}  \\
&\le&\Big(\sum_{\mathbf{u}\in\mathcal{C}(\delta)}\|D\varphi_\mathbf{u}\|^{ t_1(1-q)}\mu([\bu])^q\Big)^\alpha \Big(\sum_{\mathbf{u}\in\mathcal{C}(\delta)}\|D\varphi_\mathbf{u}\|^{ t_1(1-q)}\mu([\bu])^q \Big)^{1-\alpha},
\end{eqnarray*}
which implies that 
\begin{equation*}\begin{split}
&\overline P_\mu(\alpha t_1+(1-\alpha) t_2, q)\le\alpha \overline P_\mu(t_1, q)+(1-\alpha)\overline P_\mu( t_2, q), \quad \textit{ for $0<q<1$} \\ 
&\overline P_\mu(\alpha t_1+(1-\alpha) t_2, q)\ge\alpha \overline P_\mu(t_1, q)+(1-\alpha)\overline P_\mu( t_2, q) ,  \quad \textit{ for $q>1$}. 
\end{split}
\end{equation*}
  Hence both $\overline P_\mu(t, q)$ and  $\underline P_\mu(t, q)$  are   convex  for $0<q<1$ and  concave   for $q>1$.
\end{proof}

\end{document}
